\documentclass[12pt]{article}
\usepackage[T1]{fontenc}
\usepackage[utf8]{inputenc}
\usepackage{lmodern}
\usepackage{textcomp}
\usepackage{enumitem}
\usepackage{geometry}
\geometry{margin=1in}
\begin{document}
\begin{center}
Answer Key - Political Science - Graduate Level Assessment 1 \
\small (Answers are ordered exactly as in the question paper.)
\end{center}

\section*{Multiple Choice}
\begin{enumerate}[label=\arabic*.]
\item \textbf{Answer:} A
\\ \textit{Options:} A) It allowed the US to intensify its involvement in Vietnam; B) It illustrated the influence of public opinion on US foreign policy; C) It enhanced Congressional control over the Vietnam War; D) It curtailed US involvement in Vietnam
\\ \textit{Note:} The Gulf of Tonkin resolution granted the U.S. government broad authority to use military force in Vietnam without a formal declaration of war, thereby facilitating the escalation of American military involvement in the conflict.
\item \textbf{Answer:} D
\\ \textit{Options:} A) on nuclear weapons signed between the United States and the Soviet Union.; B) cutting conventional arms in Europe.; C) to be rejected by the U.S. Senate.; D) mandating the elimination of many long-range nuclear missiles.
\\ \textit{Note:} The Strategic Arms Reduction Treaty was groundbreaking because it specifically required the reduction and elimination of a significant number of long-range nuclear missiles, marking a pivotal shift in nuclear disarmament efforts during the Cold War.
\item \textbf{Answer:} B
\\ \textit{Options:} A) The threat posed by the Soviet Union; B) Domestic concerns of the US; C) Soviet ideology; D) None of the above
\\ \textit{Note:} Revisionists argue that the primary cause of the Cold War was rooted in the domestic concerns of the US, as American leaders prioritized internal political stability and economic interests, leading to an aggressive foreign policy aimed at containing communism.
\item \textbf{Answer:} A
\\ \textit{Options:} A) the Security Council.; B) the Chamber of Deputies.; C) the Council of Ministers.; D) the Secretariat.
\\ \textit{Note:} The Security Council holds real power within the United Nations because it is the primary body responsible for maintaining international peace and security, possessing the authority to make binding decisions, impose sanctions, and authorize military action, which other UN bodies cannot enforce.
\item \textbf{Answer:} B
\\ \textit{Options:} A) The chance for the United States to share power with other countries in the world; B) An opportunity to use to collapse of the Soviet Union to extend US power; C) An international system that didn't face any threats; D) The never-ending domination of the United States
\\ \textit{Note:} The correct answer reflects Krauthammer's assertion that the end of the Cold War and the Soviet Union's collapse created a unique geopolitical environment where the United States had the ability to shape global affairs and expand its influence unchallenged by a rival superpower.
\end{enumerate}

\section*{True / False}
\begin{enumerate}[label=\arabic*.]
\setcounter{enumi}{5}
\item \textbf{Answer:} False
\\ \textit{Options:} A) True; B) False
\\ \textit{Note:} The statement is false because, according to the latest estimates from global health organizations, there are approximately 38 million people living with HIV/AIDS worldwide, significantly exceeding the figure of 20 million.
\item \textbf{Answer:} True
\\ \textit{Options:} A) True; B) False
\\ \textit{Note:} In American politics, the Constitution grants the president the authority to recognize foreign governments, making it a key aspect of executive power in foreign affairs.
\item \textbf{Answer:} True
\\ \textit{Options:} A) True; B) False
\\ \textit{Note:} The statement is true because post-1776, American leaders and thinkers, inspired by Enlightenment ideals, viewed their nation as a model of democracy and freedom, believing it was their duty to spread these principles across the globe through expansion and influence.
\item \textbf{Answer:} True
\\ \textit{Options:} A) True; B) False
\\ \textit{Note:} Realists argue that the Cold War emerged from the power vacuum left by World War II and the Soviet Union's aggressive pursuit of hegemony, which created a clash of interests with the United States and other Western powers.
\item \textbf{Answer:} True
\\ \textit{Options:} A) True; B) False
\\ \textit{Note:} The statement is true because critics argue that the human security concept can inadvertently reinforce neo-colonial dynamics by prioritizing Western perspectives, promote global capitalism through its emphasis on market-driven solutions, and suffer from vagueness due to its broad definition that encompasses various issues without clear boundaries.
\end{enumerate}

\section*{Long Form Response}
\begin{enumerate}[label=\arabic*.]
\setcounter{enumi}{10}
\item \textbf{Answer:} The Western response to the Arab Spring highlighted several challenges to liberalism, including the rise of authoritarianism, the resurgence of nationalism, and the complexities of cultural relativism. As Western nations grappled with supporting democratic movements, they often prioritized stability over liberal ideals, leading to the backing of regimes that opposed democratic reforms. This shift not only undermined the legitimacy of Western liberalism but also contributed to a disillusionment with democratic movements in the region, as local populations perceived a lack of genuine support. Furthermore, the inconsistent application of liberal values revealed the contradictions within Western foreign policy, raising questions about the universality of liberalism in diverse cultural contexts. Collectively, these challenges have significant implications for international politics, as they shape the interactions between Western powers and emerging democracies, often favoring pragmatic approaches over principled ones.
\\ \textit{Note:} The answer correctly identifies the tension between liberal principles and the pragmatic responses of Western nations during the Arab Spring, illustrating how the prioritization of stability over democracy challenges the legitimacy of liberalism and affects international political dynamics.
\item \textbf{Answer:} The imposition of new taxes on the American colonies by Britain was largely driven by the financial burdens resulting from the war with France, specifically the Seven Years' War. This conflict significantly increased Britain's national debt and highlighted the need for revenue to cover military expenditures and colonial administration. The government viewed the American colonies as a prime source of additional income, reasoning that they had benefitted from British military protection and should therefore contribute to the costs incurred. Consequently, taxes such as the Stamp Act and the Townshend Acts were enacted, aiming to alleviate the financial strain on Britain while simultaneously asserting greater control over colonial governance. This fiscal strategy inadvertently fueled colonial resentment and the burgeoning desire for independence.
\\ \textit{Note:} The answer accurately identifies the financial strain caused by the Seven Years' War as the catalyst for Britain's decision to tax the American colonies, emphasizing the belief that the colonies, having benefited from British military efforts, should contribute to the costs incurred.
\end{enumerate}

\vfill
\noindent\textit{End of Answer Key}
\end{document}